\clearpage
\section{What the Figure 12 curves represent}\label{sec:MCpoints}
The colored curves are the ten native reconstructed-MC histograms,
normalized and plotted as connected bin densities. The local Gaussian
plus affine-pedestal fit supplies only $c$ and $\sigma_{\rm core}$;
it does not supply their line shapes. In the standardized panels, a
piecewise-linear CDF through the native 0.1~MeV histogram edges is
integrated into bins of width $\Delta u=0.02$. This is empirical
rebinning, with uniform density within each original bin, rather than
an analytic signal parameterization. The dashed black curve is a
separate candidate: the equal-mass mean of these standardized shapes.

\fig{.53}{../figures/MC_points_vs_figure12.pdf}{Direct MC-point overlays
for the outlying 60~MeV sample and a representative 160~MeV sample.
Top: full-normalized native empirical curves and counts summed into
2~MeV bins. Bottom: exactly the standardized empirical curves of
Figure 12 and direct sums into 0.2~MeV bins, both conditioned
on $|u|<2$. Horizontal bars show bin widths; vertical bars show
$\sqrt{\sum w^2}$ with the display normalization held fixed. The dashed
common shape is shown separately. Points are bin averages; the curve
must be integrated over the same bins for an exact comparison.}

The native curves agree with the MC points by construction: the
largest difference between a direct bin sum and the corresponding
CDF integral is below $10^{-9}$ candidates. This verifies the plotted
representation, not an independent goodness of fit; no fit probability
is assigned to that identity. The common-shape candidate need not
agree with each sample, and the held-out tests above assess that
separate question. The plot and all underlying counts are saved.

\clearpage
\section{The 60 MeV outlier and the meaning of a window in $u$}
At 60~MeV, $32.4\%$ of the selected histogram is inside $|u|<2$,
compared with $60.7$--$71.5\%$ at the other nine masses. Around the
current centered $\pm2.25\sigma_{\rm nominal}$ window, $5.4\%$ is
on the left, $55.1\%$ on the right and $39.5\%$ inside. Thus the
full-normalized core is small because much of this selected-candidate
distribution lies elsewhere. The core-conditioned panel removes that
fraction information and still shows some shape variation.

The 60~MeV files contain 243,668 selected rows. Their event cutflows
record 1,542,086 readout events and the same 243,668 events at the final
selection stage, a conditional retention of $15.8\%$. Other masses
give $34.5$--$61.6\%$ (Table~\ref{tab:MCunits}). This supports lower
retention from the recorded readout stage, but neither that stage nor
the unit-normalized Figure 12 establishes generated-signal acceptance.
A certified generated-event denominator and signal-daughter association
are still needed; the selected candidate distribution is not certified
as a pure signal response. ``Within the study range'' does not mean
high detector acceptance.

The existing core-centered window and a proposed fixed-$u$ window are
different because $u$ uses the fitted core width:
\begin{equation}
 |m_{ee}-c|<2.25\sigma_{\rm nominal}
 \quad\Longleftrightarrow\quad
 |u|<2.25\frac{\sigma_{\rm nominal}}{\sigma_{\rm core}}
 =2.920\text{--}3.174.
\end{equation}
Consequently $|u|<2.25$ would narrow the physical window by
$23.0$--$29.1\%$. At 60~MeV, $c=58.792$~MeV,
$\sigma_{\rm core}=1.550$~MeV and $\sigma_{\rm nominal}=2.075$~MeV:
the half-width changes from 4.669 to 3.488~MeV and the retained
selected-MC fraction falls from $39.5\%$ to $34.5\%$.

\begin{table}[H]\centering\scriptsize
\begin{tabular}{rrrrrrr}
\toprule
$m_0$ & $\sigma_{core}$ & $\sigma_{nom}$ & Old half-width ($u$) & Old MC (\%) & $\pm2.25u$ MC (\%) & Selected/readout (\%) \\
\midrule
60 & 1.550 & 2.075 & 3.012 & 39.5 & 34.5 & 15.8 \\
80 & 1.637 & 2.288 & 3.145 & 70.4 & 63.6 & 34.5 \\
100 & 1.853 & 2.570 & 3.120 & 77.6 & 71.1 & 47.5 \\
120 & 2.112 & 2.920 & 3.111 & 80.5 & 74.0 & 55.3 \\
140 & 2.392 & 3.339 & 3.140 & 81.5 & 74.6 & 59.5 \\
160 & 2.713 & 3.827 & 3.174 & 81.3 & 73.9 & 61.3 \\
180 & 3.120 & 4.383 & 3.161 & 80.7 & 72.9 & 61.6 \\
200 & 3.577 & 5.008 & 3.150 & 79.2 & 71.0 & 60.0 \\
220 & 4.233 & 5.702 & 3.031 & 77.1 & 69.4 & 57.0 \\
240 & 4.981 & 6.465 & 2.920 & 74.5 & 67.3 & 52.8 \\
\bottomrule
\end{tabular}

\caption{Widths in MeV and continuous-window probabilities from the
stored 0--400~MeV MC histograms. ``Old'' means the existing
core-centered $\pm2.25\sigma_{\rm nominal}$ window. The last column
is conditional readout-stage retention, not detector acceptance.
Actual observed-bin masks use bin centers and can differ slightly
from these continuous fractions.}\label{tab:MCunits}
\end{table}

A fixed-$u$ prescription is a coherent MC-based candidate, but the
coefficient must be declared as a new choice. It is not a change of
units that preserves the old blind region. Nor does a non-Gaussian
empirical shape inherit a Gaussian containment fraction from its fitted
core width. This clarification does not change the primary mask or
recompute observed limits and local probabilities. Testing a new mask
would require the matching GP training exclusion and signal-recovery
checks, as for the preceding width diagnostics.
